\documentclass[a4paper,11pt]{article}

\usepackage{fontspec}
\setmainfont{Carlito}
\usepackage[margin=0.6in, top=0.5in, bottom=0.5in]{geometry}
\usepackage{titlesec}
\usepackage{enumitem}
\usepackage{hyperref}
\usepackage{xcolor}
\usepackage{microtype}
\usepackage{needspace}

\hyphenpenalty=10000 \exhyphenpenalty=10000 \emergencystretch=3em

\definecolor{headernavy}{RGB}{12, 83, 166}
\definecolor{linkblue}{RGB}{0, 102, 204}
\definecolor{muted}{RGB}{96, 102, 110}

\hypersetup{colorlinks=true, linkcolor=linkblue, urlcolor=linkblue}

\titleformat{\section}
  {\color{headernavy}\normalsize\bfseries}
  {}{0em}{\MakeUppercase}[\color{headernavy}\titlerule]

\titlespacing{\section}{0pt}{8pt}{4pt}

\setlist[itemize]{noitemsep, topsep=1pt, leftmargin=1.3em, partopsep=0pt, parsep=1pt}

\pagestyle{empty}
\setlength{\parindent}{0pt}
\linespread{1.02}

% ---- building blocks: bold = organisation/project, italic = role/stack, grey = dates/places ----
% v3.5.3 (10 Nov 2026): employer (and place) on line 1, the title then the dates on line 2, so that `pdftotext -raw` reads every role as
% employer -> title -> dates -> bullets (ats_structure.role_order_problems). The dates were on line 1 before, read before the title.
\newcommand{\role}[4]{% org, location, dates, title
  \textbf{#1}{\color{muted}\,\textperiodcentered\,\textit{#2}}\\
  \textit{#4}\hfill{\color{muted}#3}\par}
\newcommand{\roleNoLoc}[3]{\textbf{#1}\\ \textit{#3}\hfill{\color{muted}#2}\par}
\newcommand{\project}[3]{% name, stack, right-hand tag
  \textbf{#1}\hfill#3\\ {\color{muted}\small\textit{#2}}\par}
\newcommand{\gh}[1]{\href{#1}{\small GitHub}}
\newcommand{\badge}[1]{{\color{headernavy}\small\textbf{#1}}}
\newcommand{\skillgroup}[2]{{\color{headernavy}\textbf{#1}}\\ #2\par\vspace{3pt}}
\newcommand{\gap}{\vspace{5pt}}
\newcommand{\num}[1]{\textbf{#1}}   % result numbers only


\begin{document}

% ================= HEADER =================
\begin{center}
  {\LARGE\bfseries Thiruvel Andagurunathan Pandian} \\[3pt]
  {\color{headernavy}\small\textbf{AI Engineer | LLM agent deployed on AWS Lambda, checked by \num{\textasciitilde{}65} tests}} \\[3pt]
  {\small
  Bristol, UK\nobreak\ \,\textbar\, \href{mailto:apthiruvel@gmail.com}{apthiruvel@gmail.com}\nobreak\ \,\textbar\, \href{tel:+447741016954}{+44 7741 016954}\nobreak\ \,\textbar\, \href{https://www.linkedin.com/in/thiruvel-a-p/}{https://www.linkedin.com/in/thiruvel-a-p/}\nobreak\ \,\textbar\, \href{https://github.com/Thiruvel-AP}{https://github.com/Thiruvel-AP}\nobreak\ \,\textbar\, \href{https://thiruvel-ap.github.io/PersonalPortfolio/}{https://thiruvel-ap.github.io/PersonalPortfolio/}}\par\vspace{4pt}
  \noindent{} Right to Work in UK \,\textbar\, Available to join Immediately
\end{center}

% ================= PROFILE =================
\needspace{5\baselineskip}\section{Profile}
AI engineer who takes LLM agents from prototype to deployment: for my AstraZeneca-affiliated MSc dissertation I deployed a tool-calling agent on AWS Lambda and checked it with a three-tier framework of \textasciitilde{}65 automated tests. At 180 Degrees Consulting, led the architecture in a five-person team that turned a startup client\textquotesingle{}s loose brief into a working agent prototype within a three-month engagement, demoed it live (5-star rating from senior leadership), then contributed to the client\textquotesingle{}s production MCP server, now in live use. Builds with AI coding tools (Claude Code, Cursor, Codex) daily. Owns data workflows end-to-end in Python and SQL, from merging structured, API and web-scraped data through feature creation, tuning, validation and deployment (Docker, GitHub Actions CI/CD, MLflow, AWS Lambda).

% ================= SKILLS =================
\needspace{5\baselineskip}\section{Skills}
\skillgroup{Agentic AI \& LLM Applications}{Tool calling and failure recovery (LangChain, LangGraph, Google ADK, ReAct agent loops; circuit breakers, fallbacks, deterministic guardrails), multi-provider LLMs (OpenAI, Anthropic Claude on AWS Bedrock, Google Gemini via LiteLLM, open-weight models via Ollama), LLM evaluation (three-tier tests, hallucination-rate analysis), agent memory (persistent storage and retrieval across sessions), RAG pipelines (retrieval, generation, embeddings and the ChromaDB vector database), MCP servers (FastMCP, TypeScript SDK), prompt engineering, AI-assisted development (Claude Code as primary tool, Cursor, OpenAI Codex)}
\skillgroup{Software Engineering \& Integrations}{TypeScript/Node.js, React, Python, JavaScript, FastAPI, Pydantic, RESTful APIs and third-party API integrations, GraphQL, PostgreSQL, SQL, MongoDB and DynamoDB (NoSQL), graph data (networkx), pytest, Git (pull requests, code review), debugging and root-cause analysis}
\skillgroup{Cloud \& Deployment}{AWS (Bedrock, Lambda, ECR, S3, SQS, EC2, IAM, Cognito), GCP (Cloud Run), Docker, Kubernetes (AWS EKS), GitHub Actions CI/CD, serverless deployment}

% ================= EXPERIENCE =================
\needspace{5\baselineskip}\section{Work Experience}
\needspace{4\baselineskip}
\role{AstraZeneca}{Bristol, UK (Remote)}{May 2026 – Sep 2026}{Data Scientist, MSc Dissertation Project}
\begin{itemize}
  \item Wrote a three-tier test framework (tool-level, API/integration, live agent) of \num{\textasciitilde{}65 tests across 8 files}, validated across the LLM provider migration from Groq to AWS Bedrock after a rate-limit evaluation.
  \item Integrated the agent into a React front end with real-time gene alias enrichment.
  \item Needed an agent that researchers could trust on a cancer research platform, including when its tools fail: built a tool-calling agent (LangChain, Claude Haiku 4.5 on AWS Bedrock) for three use cases (gene alias resolution, prompt-engineered explanations of a fixed scoring pipeline, validated dataset descriptions), with deterministic Ensembl and HGNC calls, deterministic hallucination guards, and circuit breakers with a \num{19,213-gene parquet fallback} for when the external APIs are unavailable.
  \item Built the Python data pipeline: reconciled \num{14 heterogeneous source datasets} through identifier crosswalks so the agent used trusted data, then deployed it as a Docker image on AWS Lambda (ECR, least-privilege IAM) and documented the architecture and evaluation.
\end{itemize}
\gap
\needspace{4\baselineskip}
\role{180 Degrees Consulting Bristol}{Bristol, UK (Remote)}{Sep 2025 – Dec 2025}{Data Consultant (AI Engineering \& AI System)}
\begin{itemize}
  \item Reviewed pull requests in the client\textquotesingle{}s production repository alongside the client\textquotesingle{}s CTO, the engagement\textquotesingle{}s point of contact.
  \item A product-startup client needed an agentic AI prototype turned into software their team could use: led the architecture for a five-person cross-functional team, designed hierarchical agent workflows (MCP, Exograph GraphQL, PostgreSQL for persistent storage and retrieval across agent sessions) and built the Python prototype (FastMCP, Ollama, GitHub Actions CI), then ran a live demo with client stakeholders \num{(5-star rating from senior leadership)} within the three-month engagement.
  \item Proposed the request flow: a user\textquotesingle{}s text goes through the API and a customised prompt to the MCP server, where one LLM instance writes GraphQL queries against the client\textquotesingle{}s Convex database and a second instance runs the statistical analysis before the result returns to the user.
\end{itemize}
\gap
\needspace{4\baselineskip}
\role{Freelance}{Chennai, India}{Jan 2025 – Aug 2025}{AI/ML Engineer}
\begin{itemize}
  \item Designed reproducible scikit-learn pipelines with feature engineering (temporal features, interaction terms, target encoding), model selection, cross-validation and hyperparameter tuning across classification and regression tasks.
  \item Working independently for India-based startups and individual clients, built agentic RAG workflows (LangChain, LangGraph, ChromaDB) on GCP Cloud Run, NLP pipelines (Hugging Face Transformers, Ollama) and scikit-learn pipelines, served ML models as REST APIs (FastAPI, Docker), and contributed to a team LoRA fine-tune of a Llama model.
\end{itemize}
\gap
\needspace{4\baselineskip}
\role{Avasoft}{Chennai, India}{Jan 2023 – Sep 2024}{Software Development Engineer Trainee, then Software Development Engineer}
\begin{itemize}
  \item \num{Shipped features in a production fintech (financial-services) application} as part of a team: integrated Twelve Data (real-time market feeds) and Plaid (bank linking, transaction imports) via REST APIs, with SSL pinning and AWS Cognito SSO.
  \item Delivered \num{10 reusable SwiftUI/UIKit component templates} that accelerated the team across concurrent releases, so features shipped from shared, tested building blocks.
  \item As a trainee, modelled user-profile, transaction-history and event-booking data in Core Data schemas with CRUD and session-state management, and automated a manual data-collection process with a JavaScript web-scraping pipeline that surfaced \num{market-trend signals used for pricing decisions}.
\end{itemize}

% ================= PROJECTS =================
\needspace{5\baselineskip}\section{Projects}
\needspace{4\baselineskip}
\setbox0=\hbox{\textbf{AgenticFriend: real-time voice assistant}\hspace{1em}\url{https://github.com/Thiruvel-AP/adk_multi-agent}}%
\setbox2=\hbox{{\small\textit{FastAPI, Google ADK, React, faster-whisper, PyTorch, Docker/CUDA}}\hspace{1em}{\small\url{https://github.com/Thiruvel-AP/adk_multi-agent}}}%
\ifdim\wd0>\linewidth
\ifdim\wd2>\linewidth
\textbf{AgenticFriend: real-time voice assistant}\\ {\color{muted}\small\textit{FastAPI, Google ADK, React, faster-whisper, PyTorch, Docker/CUDA}}\\ \hspace*{\fill}{\small\url{https://github.com/Thiruvel-AP/adk_multi-agent}}\par
\else
\textbf{AgenticFriend: real-time voice assistant}\\ {\color{muted}\small\textit{FastAPI, Google ADK, React, faster-whisper, PyTorch, Docker/CUDA}}\hfill{\small\url{https://github.com/Thiruvel-AP/adk_multi-agent}}\par
\fi
\else
\project{AgenticFriend: real-time voice assistant}{FastAPI, Google ADK, React, faster-whisper, PyTorch, Docker/CUDA}{\url{https://github.com/Thiruvel-AP/adk_multi-agent}}
\fi
\begin{itemize}
  \item Needed a voice assistant that answers in real time and can be interrupted: built it end to end as a full-stack application (React front end, FastAPI and Google ADK back end, full-duplex WebSocket) with ReAct-based agent orchestration, barge-in handling and self-hosted speech models on Docker/CUDA (float16 inference, model warm-up) to cut latency.
\end{itemize}
\gap
\needspace{4\baselineskip}
\setbox0=\hbox{\textbf{CoverFit: insurance benchmarking API (Bristol × UWE AI Hackathon)}\hspace{1em}\url{https://github.com/Thiruvel-AP/AIHackathon-Phoenix}}%
\setbox2=\hbox{{\small\textit{FastAPI, pandas, LiteLLM, Pydantic, pytest}}\hspace{1em}{\small\url{https://github.com/Thiruvel-AP/AIHackathon-Phoenix}}}%
\ifdim\wd0>\linewidth
\ifdim\wd2>\linewidth
\textbf{CoverFit: insurance benchmarking API (Bristol × UWE AI Hackathon)}\\ {\color{muted}\small\textit{FastAPI, pandas, LiteLLM, Pydantic, pytest}}\\ \hspace*{\fill}{\small\url{https://github.com/Thiruvel-AP/AIHackathon-Phoenix}}\par
\else
\textbf{CoverFit: insurance benchmarking API (Bristol × UWE AI Hackathon)}\\ {\color{muted}\small\textit{FastAPI, pandas, LiteLLM, Pydantic, pytest}}\hfill{\small\url{https://github.com/Thiruvel-AP/AIHackathon-Phoenix}}\par
\fi
\else
\project{CoverFit: insurance benchmarking API (Bristol × UWE AI Hackathon)}{FastAPI, pandas, LiteLLM, Pydantic, pytest}{\url{https://github.com/Thiruvel-AP/AIHackathon-Phoenix}}
\fi
\begin{itemize}
  \item An insurance sponsor needed a fast, low-cost, explainable benchmark (hackathon brief): paired a deterministic matching engine (FastAPI, pandas, Pydantic) with a single LLM call (Gemini via LiteLLM) for the rationale, covered by pytest tests; responses arrive in \num{under 3 seconds at \textasciitilde{}£0.003 per query}.
\end{itemize}

% ================= EDUCATION =================
\needspace{5\baselineskip}\section{Education}
\noindent\textit{MSc Data Science}, \textbf{University of Bristol, UK}\unskip\nobreak\hfill\penalty50\hskip1em\null\nobreak\hfill\mbox{Sep 2025 – Sep 2026}\par
\noindent{\small Dissertation with AstraZeneca affiliation (GeneTraceAI)}\par
\noindent\textit{B.E. Electrical \& Electronics Engineering}, \textbf{St. Joseph\textquotesingle{}s College of Engineering}, First Class with Distinction\unskip\nobreak\hfill\penalty50\hskip1em\null\nobreak\hfill\mbox{Aug 2019 – Apr 2023}\par

% ================= CERTIFICATIONS & ACHIEVEMENTS =================
\needspace{5\baselineskip}\section{Certifications \& Achievements}
\begin{itemize}
  \item Model Context Protocol: Advanced Topics, Anthropic, \textit{Jun 2026}
  \item Introduction to Model Context Protocol, Anthropic, \textit{Jun 2026}
  \item AWS Academy Cloud Foundations, \textit{Nov 2025}
  \item 180 Degrees Consulting Data Science Consultant Certificate, \textit{Dec 2025}
  \item Bristol PLUS Award, \textit{Mar 2026}
\end{itemize}

\end{document}